\section{Training scheme}
\label{sec:training}

\subsection{Order of fitting within a fold}

For each fold (training block, test year):
\begin{enumerate}
\item Fit the gate on the training block; freeze it; run its filter forward over the gap and the test year.
\item Fit the base with calendar-decay weights; freeze it; cache its predictions.
\item Train the experts with regime-clock weights and the frozen gate weights.
\item Predict the test year.
\end{enumerate}

\begin{whybox}[Why this order]
Each step only uses what the previous ones froze. The gate comes first because the experts' sample weights
are computed from its filtered probabilities (the regime clock) and its weights enter their likelihood.
The base comes before the experts because they learn corrections to it. Nothing fitted later can change
anything fitted earlier, which is the frozen design in procedural form.
\end{whybox}

\subsection{Windows and folds}

\dec{2026-10-05, Q16 (b)} \textbf{Expanding window, never rolling.}

\begin{whybox}
\begin{itemize}
\item \textbf{A rolling window can contain no crisis at all.} A two-year window in 2014 has no stress episode,
  so the stress expert has nothing to learn from, and by 2020 the model cannot recognise that conditions
  resemble 2008 because 2008 has left the window.
\item \textbf{Adaptation is the gate's job, not the window's.} In a regime-gated model the gate handles
  changing conditions; the base should have the longest possible sample for a low-variance estimate of the
  average relation.
\item \textbf{Old data is down-weighted, not dropped} (Section~\ref{sec:weights}), which keeps the benefits of
  recency without losing the rare episodes.
\end{itemize}
\giveup{the protocol of the closest precedent (a rolling 504-day window).}
\end{whybox}

\dec{2026-10-05} \textbf{Test period 2010--2024, refitted annually: 15 folds}, each re-fitting gate, base
and experts on the expanding block from 2000 to the year before.

\begin{whybox}
\begin{itemize}
\item \textbf{Why 2010 and not 2008}: every setting must be chosen on data before the first test year. A 2008
  start leaves only the calm years 2005--07 for tuning a regime model. A 2010 start lets the pilot see the
  2008 crisis, every training block contains two crises, and the test period still has six stress
  episodes, including COVID, the sharpest transition test. 2008 itself can still be tested out of sample
  by leave-one-episode-out.
\item \textbf{Why annual}: the filter already updates the regime every day, so the model reacts daily;
  refits only update slow parameters. Quarterly refits would cost four times as much for little gain, and
  Gu, Kelly \& Xiu also refit annually with an expanding window.
\end{itemize}
\end{whybox}

A \textbf{5-day purge} separates training and test blocks (the gap's labels are dropped; the gate's filter
still sees its returns). \open{Q16 (c)} \textbf{Leave-one-episode-out}, as headline or robustness.

\begin{whybox}[Why the purge]
A training row from four days before the test year has a target that ends inside the test year. Keeping it
would let test-period returns leak into training. Removing $h$ days removes every such overlap.
\end{whybox}

\subsection{Sample weights: one memory per component}
\label{sec:weights}

\dec{2026-10-05, Q16 (d)(e)}

{\small
\begin{tabularx}{\textwidth}{@{}l X X X@{}}
\toprule
Component & Estimates & Main risk & Memory\\
\midrule
Base & the average cross-sectional mapping & bias: structural drift shows here & calendar exponential
decay, long half-life\\
Experts & regime-specific mappings & variance for rare regimes & regime-clock decay\\
Gate & regime dynamics of the market & both & regime-clock forgetting\\
\bottomrule
\end{tabularx}}

\begin{whybox}[Why weight at all, and why differently per component]
This is a bias--variance trade-off. Old data adds \emph{bias} because markets change: anomalies weaken after
publication (McLean \& Pontiff 2016) and as markets become more liquid (Chordia, Subrahmanyam \& Tong 2014).
Dropping old data adds \emph{variance} in a setting where the signal is tiny, and the optimal weights after a
structural break are smaller but not zero (Pesaran \& Timmermann 2007). Exponential down-weighting sits
between the expanding window (no decay) and a rolling window (a hard cut-off). The three components face
the trade-off differently, so each gets its own memory.
\end{whybox}

\paragraph{Calendar decay (base).} A row $a$ trading days old gets weight $w_a=\rho^{a}$ with $\rho=2^{-1/H}$.
The Kish effective sample size then tends to $(1+\rho)/(1-\rho)$, about 2.9 half-lives, however long the
history.

\begin{whybox}
The base estimates the average relation across all regimes, which is where slow structural drift shows
most, and every day informs it, so ordinary recency weighting fits. The half-life is long because the
signal is weak and the effective sample shrinks fast as the half-life shortens.
\end{whybox}

\paragraph{Regime-clock decay (experts; Tom's construction).} A row's age is the amount of \emph{later
experience of the same regime}, counted with the gate's filtered probabilities:
\begin{equation}
w_s(t)=\sum_{k=1}^{K}\xi_{s\mid s}(k)\,\rho^{\,n_k(s,t)},\qquad
n_k(s,t)=\sum_{u=s+1}^{t}\xi_{u\mid u}(k),\qquad \rho=2^{-1/H}.
\end{equation}
A regime occupying a share $p$ of days has a calendar half-life of about $H/p$: with $H=500$ regime-days,
about 2.8 years for a calm regime (70\% of days), 6.6 years for stress (30\%), decades for a rare crisis
state. \open{minor} mixture form (default) or a per-expert clock.

\begin{whybox}
\begin{itemize}
\item \textbf{Calendar decay would erase the crises}: with a two-year half-life, 2008 data carries about 1.6\%
  weight by 2020. The crisis regime has the least data and almost all of it is old, so calendar decay
  starves exactly the expert that needs history most.
\item \textbf{Regime-clock decay resolves the conflict between ``markets drift'' and ``regimes recur''}: in a
  long calm spell, calm data turns over quickly (drift is tracked) while the last crisis keeps its weight
  because no newer crisis experience has replaced it. After a cluster of crises, older crisis data fades.
\item \textbf{Nothing is hand-labelled}: the gate's own filtered probabilities do the work, with no look-ahead
  and no choice of ``important periods'' that could be snooped.
\end{itemize}
The closest precedent for weighting history by regime rather than calendar time is Mulliner, Harvey, Xia,
Fang \& Van Hemert (2025). Measuring age in regime time with filtered probabilities is this study's own
construction.
\end{whybox}

In the likelihood each row's term is multiplied by its weight (a weighted composite likelihood): scores
stay unbiased, so the weights change efficiency, not consistency.

\subsection{The pilot: every setting chosen once}

\dec{2026-10-05} A pilot on three folds with \textbf{validation years 2007, 2008 and 2009}, which never touches
2010 or later, chooses once: the half-lives (from a small grid always including equal weights) and the gate
memory; learning rate, weight decay and batch size; the training length as a fixed number of steps.
Settings other than the gate's memory are chosen by a \textbf{regime-balanced validation loss}, each row
reweighted by
\begin{equation}
v(s)=\sum_{k}\frac{\bar w_k(s)}{M_k},
\end{equation}
where $\bar w_k(s)$ is the window-average gate weight and $M_k$ its mean over the validation block.

\begin{whybox}
\begin{itemize}
\item \textbf{Choose once, before the test period}: tuning in every fold or on test data would let test results
  shape the model (Q15). One lock file, written before any test year is touched, is the pre-registration.
\item \textbf{Why 2007--2009}: the validation years include the calm run-up, the crisis and the recovery, so
  settings are chosen on a full cycle.
\item \textbf{Why regime-balanced}: about 70\% of days are calm, so an unweighted validation loss would choose
  settings that serve calm markets and ignore stress, the very thing the gate exists for. Balancing makes
  every regime count equally.
\item \textbf{Equal weights always in the grid}, so ``decay does not help'' is a possible, reportable outcome, and
  the undecayed protocol of Gu, Kelly \& Xiu is always available as a reference.
\end{itemize}
\end{whybox}

\subsection{Main study: frozen settings, full blocks}

\dec{2026-10-05}
\begin{itemize}
\item \textbf{No held-out validation tail and no per-fold early stopping}; fallback to regime-balanced early
  stopping if the pilot's best step count drifts across its folds.
\item \textbf{Train on every day} under a fixed step budget.
\item \textbf{Fresh initialisation every fold}; no warm starts (fallback: shrink-and-perturb).
\item \textbf{Paired seeds}: 10 per grid cell, 30 for the primary cell, final count from the pilot.
\item \textbf{One process per performance core}, large batches, features built once, base cached.
\end{itemize}

\begin{whybox}
\begin{itemize}
\item \textbf{No validation tail}: holding out the last 20\% of each block would withhold the most recent, most
  heavily weighted data from training. Decay keeps the effective sample roughly stable as the window grows,
  so a step budget chosen in the pilot should transfer; the pilot checks this.
\item \textbf{Every day}: under a fixed step budget, thinning to every fifth day saves no compute; overlapping
  rows still add information because the price features change daily, and overlap only affects standard
  errors, which are corrected.
\item \textbf{Fresh starts}: warm-started networks generalise worse (Ash \& Adams 2020); a warm start would also
  make corrections relative to last year's base, breaking the zero-initialised residual design, and would
  tie the folds together.
\item \textbf{Paired seeds}: the quantity of interest is the \emph{difference} between gates. Running every arm on
  the same seeds makes much of the training noise common to both arms, so it cancels in the difference
  (common random numbers). The seed count is set by the measured noise, not guessed.
\item \textbf{Process parallelism}: measured on the laptop, one thread per run with ten runs in parallel gives
  about seven times the throughput; extra threads inside one run give nothing.
\end{itemize}
\end{whybox}
